\documentclass[10pt,a4paper]{amsart}
\usepackage[margin=19mm]{geometry}
\usepackage{amsmath,amssymb,mathtools}
\usepackage[scr=rsfs,cal=euler]{mathalfa}
\usepackage{xfrac}
\usepackage[unicode,hidelinks,bookmarks=false]{hyperref}
\newcommand{\ie}{i.e.\ }
\newcommand{\cf}{cf.\ }
\newcommand{\resp}{respectively\ }
\newcommand{\Subsection}{\subsection}
\providecommand{\nobreakdash}{\nobreak\hbox{-}}
\setlength{\emergencystretch}{2em}
\multlinegap0pt
\usepackage{lmodern}
\usepackage{mathtools}
\usepackage{booktabs,tabularx,array,multirow}
\usepackage{enumitem}
\usepackage{xurl}
\usepackage{fancyhdr}
\usepackage{cleveref}
\setlength{\parindent}{1.5em}
\setlength{\parskip}{0.25em}
\setlength{\emergencystretch}{2em}
\allowdisplaybreaks[2]
\numberwithin{equation}{section}
\newtheorem{theorem}{Theorem}[section]
\newtheorem{proposition}[theorem]{Proposition}
\newtheorem{lemma}[theorem]{Lemma}
\newtheorem{corollary}[theorem]{Corollary}
\newtheorem{definition}[theorem]{Definition}
\newtheorem{example}[theorem]{Example}
\newtheorem{remark}[theorem]{Remark}
\newcommand{\I}{\mathbb I}
\newcommand{\Sier}{\mathbb S}
\newcommand{\LFTS}{\mathbf{LFTS}}
\newcommand{\const}[1]{\underline{#1}}
\newcommand{\Fin}{\operatorname{Fin}}
\newcommand{\ev}{\operatorname{ev}}
\newcommand{\sub}{\operatorname{sub}}
\newcommand{\Hom}{\operatorname{Hom}}
\newcommand{\id}{\operatorname{id}}
\begin{document}
\title[Exponentiable Objects and Function spaces in Lowen Fuzzy Top]{Exponentiable Objects and Function spaces in Lowen Fuzzy Topological Spaces}
\author{Yongming Li}
\date{}
\maketitle
\noindent\emph{Source and licence.} Exponentiable Objects and Function spaces in Lowen Fuzzy Topological Spaces. Version arXiv:2609.23281v1 (CC BY 4.0). This material is distributed under the Creative Commons Attribution 4.0 International licence, http://creativecommons.org/licenses/by/4.0/, as recorded in the arXiv OAI licence metadata for this exact version and shown on the abstract page https://arxiv.org/abs/2609.23281v1. Full rights notice: licenses/arxiv-2609.23281-CC-BY-4.0.txt.\par\smallskip
\noindent\textit{Editorial scope.} Counted unit: the original \texttt{theorem} environment \texttt{thm:main-en} at source lines 478--497 together with its complete proof at lines 499--536; the delivered document keeps source lines 227--537 so that every referenced label and every citation resolves inside it. Native editable source is the paper's own concatenated TeX; the AMS wrapper is editorial formatting only. No publisher class, upstream script or unrelated artwork is redistributed.
\begin{lemma}[Rectangles in products]\label{lem:rect-en}
For Lowen spaces \(Z,X\), a map \(w:Z\times X\to\I\) is open in the
categorical product if and only if
\begin{equation}\label{eq:rect-en}
 w(z,x)=\bigvee_{j\in J}\bigl(a_j(z)\wedge b_j(x)\bigr),
 \qquad a_j\in\tau_Z,\quad b_j\in\tau_X.
\end{equation}
\end{lemma}

\begin{proof}
The product topology is generated by the inverse images of fuzzy opens
under the two projections.  A finite meet of such generators combines
into one factor from \(Z\) and one from \(X\); arbitrary joins then give
\eqref{eq:rect-en}.  The reverse inclusion follows from closure of the
product topology.
\end{proof}

\subsection{The Lowen fuzzy Sierpiński object}

Let \(\delta_{\Sier}\) be the least Lowen topology on \(\I\) generated
by \(\id_{\I}\), and put \(\Sier=(\I,\delta_{\Sier})\).  Early fuzzy
Sierpiński constructions and their categorical role appear in
~\cite{Srivastava1984-en,SrivastavaSrivastava1986-en,
LowenSrivastava1988-en}; the corresponding stratified \(Q\)-topological
form is discussed in~\cite{TiwariSrivastava2022-en}.

\begin{proposition}[Classification of fuzzy opens]\label{prop:class-en}
For every Lowen space \(X\) and every map \(u:X\to\I\),
\[
 u\in\tau_X \quad\Longleftrightarrow\quad u:X\to\Sier
 \text{ is continuous}.
\]
Consequently, naturally in \(X\),
\begin{equation}\label{eq:hom-en}
 C(X,\Sier)=\Hom_{\LFTS}(X,\Sier)\cong\tau_X.
\end{equation}
\end{proposition}

\begin{proof}
If \(u\) is continuous, then \(u=\id_\I\circ u\) is fuzzy open.  The
converse follows because every member of \(\delta_{\Sier}\) is obtained
from \(\id_\I\) and constants by arbitrary joins and finite meets.
\end{proof}

\begin{lemma}[Fuzzy opens of \(\Sier^n\)]\label{lem:sier-n-en}
For \(n\geq1\), a map \(p:\I^n\to\I\) is fuzzy open in the product
\(\Sier^n\) if and only if there are constants
\(c_A\in\I\), \(A\subseteq\{1,\ldots,n\}\), such that
\begin{equation}\label{eq:sugeno-en}
 p(r_1,\ldots,r_n)=
 \bigvee_{A\subseteq\{1,\ldots,n\}}
 \left(c_A\wedge\bigwedge_{i\in A}r_i\right),
\end{equation}
where the empty meet is \(1\).  In particular, the unary fuzzy opens
are exactly
\begin{equation}\label{eq:clamp-en}
 p(r)=a\vee(r\wedge b),\qquad 0\leq a\leq b\leq1.
\end{equation}
\end{lemma}

\begin{proof}
The product is generated by the coordinate maps and the constants.
Every finite meet of generators has the form
\(c\wedge\bigwedge_{i\in A}r_i\).  Terms with the same finite index set
can be joined by taking the supremum of their coefficients, which gives
\eqref{eq:sugeno-en}.  The converse is immediate.
\end{proof}

\section{Function spaces, splitting, and exponentiability}

Let \(C(X,Y)\) denote the set of continuous maps from \(X\) to \(Y\).
A Lowen topology \(\eta\) on this set is \emph{splitting} if every
continuous \(H:Z\times X\to Y\) has continuous transpose
\(\widehat H:Z\to(C(X,Y),\eta)\).  It is \emph{conjoining} if the
evaluation map
\[
 \ev:C(X,Y)\times X\to Y,\qquad \ev(f,x)=f(x),
\]
is continuous.

\begin{lemma}[Exponential-law criterion]\label{lem:exp-law-en}
An object \(X\) is exponentiable in \(\LFTS\) if and only if, for every
Lowen space \(Y\), the set \(C(X,Y)\) admits a topology that is both
splitting and conjoining.  That function space is then \(Y^X\).
\end{lemma}

\begin{proof}
The forward implication is the exponential adjunction and its counit.
Conversely, transposition and evaluation are inverse on underlying sets
and are natural once splitting and conjoining hold.  The underlying set
of an exponential (denoted by $(U(Y^X)$ here) is indeed \(C(X,Y)\): every point selection from the
one-point terminal Lowen space is continuous, so
\(U(Y^X)\cong\Hom(1,Y^X)\cong\Hom(X,Y)\).
\end{proof}

The preceding criterion is the concrete form, in \(\LFTS\), of the
proper/admissible and splitting/conjoining principles of Schwarz and
Alderton~\cite{Schwarz1983-en,Alderton1988-en}.  Alderton subsequently
applied this framework to fuzzy function spaces~\cite{Alderton1989-en},
and J\"ager compared it with continuous-convergence and compact-open
constructions~\cite{Jager1999-en}.  What remains here is to calculate
the largest splitting topology and the conjoining condition explicitly.

Splitting topologies are closed under taking the Lowen topology generated
by their union.  Hence \(C(X,Y)\) has a largest splitting topology.  We
now compute it for \(Y=\Sier\).

\section{Tier-compatible Scott weights}

Regard \(\tau_X\) as a complete lattice in the pointwise order.  Directed
sets are nonempty.

\begin{definition}[Scott weight and finite-tier compatibility]
A map \(\Phi:\tau_X\to\I\) is Scott-continuous if it is monotone and
\[
 \Phi(\bigvee D)=\bigvee_{\mu\in D}\Phi(\mu)
\]
for every directed \(D\subseteq\tau_X\).  Such a weight satisfies
\emph{finite-tier compatibility} (FTC) if, for every \(n\geq1\) and
\(\beta_1,\ldots,\beta_n\in\tau_X\), the map
\begin{equation}\label{eq:ftc-en}
 P_{\Phi,\boldsymbol\beta}(r_1,\ldots,r_n)
 =\Phi\left(\bigvee_{i=1}^n
      (\const{r_i}\wedge\beta_i)\right)
\end{equation}
is fuzzy open in \(\Sier^n\).  Let \(\Sigma_X\) be the set of all
Scott-continuous weights satisfying FTC.
\end{definition}

By Lemma~\ref{lem:sier-n-en}, FTC says exactly that every finite-tier
test \eqref{eq:ftc-en} is a finite lattice polynomial of the form
\eqref{eq:sugeno-en}; it does not impose invariance under arbitrary
affine regradings of \([0,1]\).

\begin{proposition}\label{prop:sigma-en}
\(\Sigma_X\subseteq\I^{\tau_X}\) is a Lowen fuzzy topology.
\end{proposition}

\begin{proof}
Constant weights satisfy the conditions.  Arbitrary joins commute with
directed joins and with all finite-tier tests.  If \(\Phi,\Psi\) satisfy
the conditions, directedness gives
\[
 \Phi\left(\bigvee D\right)\wedge \Psi\left(\bigvee D\right)
 =\bigvee_{\mu\in D}(\Phi(\mu)\wedge\Psi(\mu));
\]
hence \(\Phi\wedge\Psi\) is Scott-continuous, and its FTC tests are
finite meets of fuzzy opens.
\end{proof}

\begin{theorem}[Product stability]\label{thm:prod-stable-en}
For a map \(\Phi:\tau_X\to\I\), the following are equivalent:
\begin{enumerate}[label=(\roman*)]
 \item \(\Phi\in\Sigma_X\);
 \item for every Lowen space \(Z\) and every
 \(w\in\tau_{Z\times X}\), the map
 \[
  z\longmapsto\Phi(w_z),\qquad w_z(x)=w(z,x),
 \]
 belongs to \(\tau_Z\).
\end{enumerate}
\end{theorem}

\begin{proof}
Assume (i) and write \(w\) as in \eqref{eq:rect-en}.  For a nonempty
finite \(F\subseteq J\), put
\[
 w_{F,z}=\bigvee_{j\in F}
       (\const{a_j(z)}\wedge b_j).
\]
These finite partial joins form a directed family with supremum \(w_z\).
Scott continuity and FTC give
\[
 \Phi(w_z)=\bigvee_{F\in\Fin_+(J)}\Phi(w_{F,z})\in\tau_Z.
\]
Indeed, for \(F=\{j_1,\ldots,j_n\}\), the corresponding summand is the
composite of \(Z\to\Sier^n\),
\(z\mapsto(a_{j_1}(z),\ldots,a_{j_n}(z))\), with the FTC test of
\(\Phi\) at \((b_{j_1},\ldots,b_{j_n})\).

Conversely, taking \(Z=\Sier^n\) and
\(w(\boldsymbol r,x)=\bigvee_i(r_i\wedge\beta_i(x))\) proves FTC.  To
prove Scott continuity, let \(D\subseteq\tau_X\) be directed and form
\(Z_D=D\sqcup\{\infty\}\).  Generate a Lowen topology on \(Z_D\) by
the constants and
\[
 u_d(e)=1\ (d\leq e\text{ or }e=\infty),
 \qquad u_d(e)=0\text{ otherwise}.
\]
Every fuzzy open \(q\) in this topology is monotone on \(D\) and
satisfies \(q(\infty)=\bigvee_{e\in D}q(e)\).  The product-open map
\[
 W(z,x)=\bigvee_{d\in D}(u_d(z)\wedge d(x))
\]
has sections \(W_e=e\) and \(W_\infty=\bigvee D\).  Applying (ii) to
\(W\) yields
\(\Phi(\bigvee D)=\bigvee_{e\in D}\Phi(e)\).  Applying this identity to
the two-element directed set \(\{\mu,\nu\}\), \(\mu\leq\nu\), also
gives monotonicity.
\end{proof}

\begin{theorem}[Largest splitting topology]\label{thm:max-split-en}
Under \(C(X,\Sier)=\tau_X\), the Lowen topology \(\Sigma_X\) is the
largest splitting topology on \(C(X,\Sier)\).
\end{theorem}

\begin{proof}
Proposition~\ref{prop:sigma-en} and the forward implication of
Theorem~\ref{thm:prod-stable-en} show that \(\Sigma_X\) is splitting.
If \(\eta\) is any splitting topology and \(\Phi\in\eta\), apply
splitting to each fuzzy open \(w:Z\times X\to\Sier\).  The pullback of
\(\Phi\) along its transpose is \(z\mapsto\Phi(w_z)\); the reverse
implication of Theorem~\ref{thm:prod-stable-en} gives
\(\Phi\in\Sigma_X\).  Thus \(\eta\subseteq\Sigma_X\).
\end{proof}

The following example shows that Scott continuity of a weight is not FTC.

\begin{example}[A one-point weight test]
Let \(X=1\).  Then \(\tau_X\cong\I\).  The weight
\(\Phi(t)=t^2\) is monotone and preserves directed suprema, hence is
Scott-continuous.  For \(n=1\), \(\beta_1=1\), its FTC test is
\(P(r)=r^2\).  If \(r^2=a\vee(r\wedge b)\), substitution of \(0\) and
\(1\) forces \(a=0,b=1\), which would give \(r^2=r\), a contradiction.
Thus \(\Phi\notin\Sigma_X\).
\end{example}

This is a weight-level test, not a space-level counterexample: the
one-point object is exponentiable.  Indeed, the identity weight and
\(\lambda=1\) satisfy SKA below.  The example isolates the finite-tier
condition missing from an unrestricted Scott function-space topology.

\section{The stratified Scott--core theorem}

\begin{definition}[Compatible core and SKA]\label{def:ska-en}
For \(\lambda\in\tau_X\) and \(\Phi\in\Sigma_X\), write
\begin{equation}\label{eq:triangle-en}
 \lambda\triangleleft\Phi
 \quad\Longleftrightarrow\quad
 \const{\Phi(\nu)}\wedge\lambda\leq\nu
 \quad\text{for all }\nu\in\tau_X.
\end{equation}
The space \(X\) satisfies stratified Scott--core approximation (SKA) if
\begin{equation}\label{eq:ska-en}
  \mu=\bigvee_{\substack{\Phi\in\Sigma_X,\ \lambda\in\tau_X\\
                         \lambda\triangleleft\Phi}}
       (\const{\Phi(\mu)}\wedge\lambda)
 \qquad(\mu\in\tau_X).
\end{equation}
\end{definition}


\begin{theorem}[Characterization of exponentiable objects]
\label{thm:main-en}
For a Lowen fuzzy topological space \(X\), the following are equivalent:
\begin{enumerate}[label=(\roman*)]
 \item \(X\) is exponentiable in \(\LFTS\);
 \item the evaluation map
 \[
 e:(\tau_X,\Sigma_X)\times X\to\Sier,
 \qquad e(\mu,x)=\mu(x),
 \]
 is continuous;
 \item \(X\) satisfies SKA.
\end{enumerate}
When these conditions hold, \(Y^X\) has underlying set \(C(X,Y)\),
with the Lowen topology generated by the constants and
\begin{equation}\label{eq:exp-gen-en}
 [\Phi,v](f)=\Phi(v\circ f),
 \qquad \Phi\in\Sigma_X,\quad v\in\tau_Y.
\end{equation}
\end{theorem}

\begin{proof}
Assume (i).  The topology of \(\Sier^X\), transported to \(\tau_X\),
is splitting and its evaluation is continuous.  It is contained in the
largest splitting topology \(\Sigma_X\).  Refining the function-space
factor preserves the openness of the evaluation pullbacks, proving (ii).

Assume (ii).  By Lemma~\ref{lem:rect-en}, the fuzzy open underlying
evaluation admits a rectangle decomposition
\[
 e(\mu,x)=\bigvee_{j\in J}
       (\Phi_j(\mu)\wedge\lambda_j(x)),
 \qquad \Phi_j\in\Sigma_X,\quad\lambda_j\in\tau_X.
\]
Every rectangle is bounded above by \(e\); hence
\(\const{\Phi_j(\nu)}\wedge\lambda_j\leq\nu\) for all \(\nu\), so
\(\lambda_j\triangleleft\Phi_j\).  Fixing \(\mu\) gives both inequalities
in \eqref{eq:ska-en} and proves (iii).

Assume (iii), fix \(Y\), and give \(C(X,Y)\) the topology generated by
\eqref{eq:exp-gen-en}.  If \(H:Z\times X\to Y\) is continuous, then for
\(v\in\tau_Y\) the map \(w(z,x)=v(H(z,x))\) is product-open.  Product
stability gives
\[
 [\Phi,v](\widehat H(z))=\Phi(w_z)\in\tau_Z,
\]
so the topology is splitting.  For evaluation, put \(\mu=v\circ f\).
SKA gives
\begin{align*}
 v(\ev(f,x))
 &=\mu(x)\\
 &=\bigvee_{\lambda\triangleleft\Phi}
   (\Phi(\mu)\wedge\lambda(x))\\
 &=\bigvee_{\lambda\triangleleft\Phi}
   ([\Phi,v](f)\wedge\lambda(x)),
\end{align*}
a join of open rectangles.  Evaluation is continuous, and
Lemma~\ref{lem:exp-law-en} supplies the full exponential adjunction.
\end{proof}

\begin{thebibliography}{99}
\bibitem[Alderton1988-en]{Alderton1988-en} I. W. Alderton, Splitting and conjoining objects in monotopological categories, \emph{Topology and its Applications} \textbf{29} (1988), 223--235, doi:10.1016/0166-8641(88)90022-3.
\bibitem[Alderton1989-en]{Alderton1989-en} I. W. Alderton, Function spaces in fuzzy topology, \emph{Fuzzy Sets and Systems} \textbf{32} (1989), 115--124, doi:10.1016/0165-0114(89)90092-4.
\bibitem[Jager1999-en]{Jager1999-en} G. J\"ager, On fuzzy function spaces, \emph{International Journal of Mathematics and Mathematical Sciences} \textbf{22} (1999), 727--737, doi:10.1155/S0161171299227275.
\bibitem[LowenSrivastava1988-en]{LowenSrivastava1988-en} R. Lowen and A. K. Srivastava, Sierpinski objects in subcategories of FTS, \emph{Quaestiones Mathematicae} \textbf{11} (1988), 181--193.
\bibitem[Schwarz1983-en]{Schwarz1983-en} F. Schwarz, Powers and exponential objects in initially structured categories and applications to categories of limit spaces, \emph{Quaestiones Mathematicae} \textbf{6} (1983), 227--254, doi:10.1080/16073606.1983.9632302.
\bibitem[Srivastava1984-en]{Srivastava1984-en} A. K. Srivastava, Fuzzy topology and Sierpinski fuzzy space, \emph{Journal of Mathematical Analysis and Applications} \textbf{103} (1984), 103--117, doi:10.1016/0022-247X(84)90160-4.
\bibitem[SrivastavaSrivastava1986-en]{SrivastavaSrivastava1986-en} R. Srivastava and A. K. Srivastava, The Sierpinski object in fuzzy topology, in \emph{Modern Analysis and Its Applications}, Prentice-Hall of India, New Delhi, 1986, pp.~17--25.
\bibitem[TiwariSrivastava2022-en]{TiwariSrivastava2022-en} H. Tiwari and R. Srivastava, On coreflective hulls in Str-\(Q\)-TOP, \emph{Soft Computing} \textbf{26} (2022), 527--534, doi:10.1007/s00500-021-06537-z.
\end{thebibliography}
\end{document}
